% GENERATED FILE. Edit canonical lessons, then run npm run generate:books.
% canonical-source-hash: fnv1a64:f3399713b51f01ac
% canonical-lessons: ML-C486-embasi, ML-C486-vakuppu, ML-C486-sakha, ML-C486-pravrttisamayam, ML-C486-pravrttidinam

\chapter{Embassies and Offices}
\label{ch:ml-embassies-and-offices}
\hlchaptermodality{486}

\begin{chapteropening}
\textbf{By the end of this chapter:} \emph{I can say an embassy, a department, a branch office, office hours and a working day.}

Complete the last lesson of chapter 486: I can say an embassy, a department, a branch office, office hours and a working day.
\end{chapteropening}

\section[\texorpdfstring{\ml{എംബസി} (embasi)}{എംബസി (embasi)}]{\ml{എംബസി} (embasi) — an embassy}

\label{lesson:ML-C486-embasi}



\begin{quote}
Before the new one: say the Malayalam for a waiting room, then the Malayalam for a roadside food stall.
\end{quote}

\subsection*{You'll want to know: \ml{എംബസി}}
\textbf{\ml{എംബസി}} — \emph{embasi} — \textquotedblleft{}an embassy\textquotedblright{}.

\begin{quote}
\ml{വിസയ്ക്കായി} \ml{അവൾ} \ml{എംബസിയിൽ} \ml{പോയി}.
\end{quote}

\begin{quote}
\emph{visaykkāyi avaḷ embasiyil pōyi.}
\end{quote}

\begin{quote}
She went to the embassy for a visa.
\end{quote}

Each country has one embassy in the capital, and a consulate in some other cities.

\begin{cousinweb}
\textbf{\ml{എംബസി}} (\emph{embasi}) was borrowed from English \emph{embassy}, the office of a foreign government in the capital.
\end{cousinweb}

\begin{grammarlens}[title={What you've built}]
One more word for getting around.
\end{grammarlens}

\subsection*{Your turn}
\begin{itemize}
\raggedright
  \item \emph{Say it:} \emph{embasi}
  \item \emph{Say it:} \emph{embasi}, once more
  \item \emph{Say it:} \emph{visaykkāyi avaḷ embasiyil pōyi.} — She went to the embassy for a visa.
  \item \emph{Recall:} say the Malayalam for a waiting room, then the Malayalam for a roadside food stall, then say \emph{embasi} again
\end{itemize}

\begin{tcolorbox}[breakable,colback=teal!4,colframe=teal!35!black,title={Before you move on}]
What does \ml{എംബസി} mean? (\textquotedblleft{}An embassy\textquotedblright{}.) Say it once more.
\end{tcolorbox}
\section[\texorpdfstring{\ml{വകുപ്പ്} (vakuppŭ)}{വകുപ്പ് (vakuppŭ)}]{\ml{വകുപ്പ്} (vakuppŭ) — a department}

\label{lesson:ML-C486-vakuppu}



\begin{quote}
Before the new one: say the Malayalam for a roadside food stall, then the Malayalam for an embassy.
\end{quote}

\subsection*{You'll want to know: \ml{വകുപ്പ്}}
\textbf{\ml{വകുപ്പ്}} — \emph{vakuppŭ} — \textquotedblleft{}a department\textquotedblright{}.

\begin{quote}
\ml{റെയിൽവേ} \ml{ഒരു} \ml{വലിയ} \ml{സർക്കാർ} \ml{വകുപ്പാണ്}.
\end{quote}

\begin{quote}
\emph{ṟeyilvē oru valiya sarkkār vakuppāṇŭ.}
\end{quote}

\begin{quote}
The railway is a big government department.
\end{quote}

It also means a section of a law: \ml{നിയമത്തിന്റെ} \ml{ഒരു} \ml{വകുപ്പ്}.

\begin{cousinweb}
\textbf{\ml{വകുപ്പ്}} (\emph{vakuppŭ}) is native Dravidian, from the root \ml{വക}, to divide; DEDR gives the Malayalam noun as a section or kind (DEDR 5202).
\end{cousinweb}

\begin{grammarlens}[title={What you've built}]
One more word for getting around.
\end{grammarlens}

\subsection*{Your turn}
\begin{itemize}
\raggedright
  \item \emph{Say it:} \emph{vakuppŭ}
  \item \emph{Say it:} \emph{vakuppŭ}, once more
  \item \emph{Say it:} \emph{ṟeyilvē oru valiya sarkkār vakuppāṇŭ.} — The railway is a big government department.
  \item \emph{Recall:} say the Malayalam for a roadside food stall, then the Malayalam for an embassy, then say \emph{vakuppŭ} again
\end{itemize}

\begin{tcolorbox}[breakable,colback=teal!4,colframe=teal!35!black,title={Before you move on}]
What does \ml{വകുപ്പ്} mean? (\textquotedblleft{}A department\textquotedblright{}.) Say it once more.
\end{tcolorbox}
\section[\texorpdfstring{\ml{ശാഖ} (śākha)}{ശാഖ (śākha)}]{\ml{ശാഖ} (śākha) — a branch office; a branch}

\label{lesson:ML-C486-sakha}



\begin{quote}
Before the new one: say the Malayalam for an embassy, then the Malayalam for a department.
\end{quote}

\subsection*{You'll want to know: \ml{ശാഖ}}
\textbf{\ml{ശാഖ}} — \emph{śākha} — \textquotedblleft{}a branch office; a branch\textquotedblright{}.

\begin{quote}
\ml{ആ} \ml{ബാങ്കിന്} \ml{ഞങ്ങളുടെ} \ml{ഗ്രാമത്തിൽ} \ml{ഒരു} \ml{ശാഖ} \ml{ഉണ്ട്}.
\end{quote}

\begin{quote}
\emph{ā bāṅkinŭ ñaṅṅaḷuṭe grāmattil oru śākha uṇṭŭ.}
\end{quote}

\begin{quote}
That bank has a branch in our village.
\end{quote}

In formal Malayalam it is also the branch of a tree; in speech that is a \ml{കൊമ്പ്}.

\begin{cousinweb}
\textbf{\ml{ശാഖ}} (\emph{śākha}) is borrowed from Sanskrit \emph{śākhā}, \textquotedblleft{}a branch\textquotedblright{} of a tree, and so a part of a larger body (Monier-Williams).
\end{cousinweb}

\begin{grammarlens}[title={What you've built}]
One more word for getting around.
\end{grammarlens}

\subsection*{Your turn}
\begin{itemize}
\raggedright
  \item \emph{Say it:} \emph{śākha}
  \item \emph{Say it:} \emph{śākha}, once more
  \item \emph{Say it:} \emph{ā bāṅkinŭ ñaṅṅaḷuṭe grāmattil oru śākha uṇṭŭ.} — That bank has a branch in our village.
  \item \emph{Recall:} say the Malayalam for an embassy, then the Malayalam for a department, then say \emph{śākha} again
\end{itemize}

\begin{tcolorbox}[breakable,colback=teal!4,colframe=teal!35!black,title={Before you move on}]
What does \ml{ശാഖ} mean? (\textquotedblleft{}A branch office; a branch\textquotedblright{}.) Say it once more.
\end{tcolorbox}
\section[\texorpdfstring{\ml{പ്രവൃത്തിസമയം} (pravṛttisamayaṁ)}{പ്രവൃത്തിസമയം (pravṛttisamayaṁ)}]{\ml{പ്രവൃത്തിസമയം} (pravṛttisamayaṁ) — office hours; working hours}

\label{lesson:ML-C486-pravrttisamayam}



\begin{quote}
Before the new one: say the Malayalam for a department, then the Malayalam for a branch office.
\end{quote}

\subsection*{You'll want to know: \ml{പ്രവൃത്തിസമയം}}
\textbf{\ml{പ്രവൃത്തിസമയം}} — \emph{pravṛttisamayaṁ} — \textquotedblleft{}office hours; working hours\textquotedblright{}.

\begin{quote}
\ml{ബാങ്കിന്റെ} \ml{പ്രവൃത്തിസമയം} \ml{രാവിലെ} \ml{പത്ത്} \ml{മുതൽ} \ml{നാല്} \ml{വരെയാണ്}.
\end{quote}

\begin{quote}
\emph{bāṅkinṟe pravṛttisamayaṁ rāvile pattŭ mutal nālŭ vareyāṇŭ.}
\end{quote}

\begin{quote}
The bank's working hours are from ten in the morning to four.
\end{quote}

\ml{പ്രവൃത്തി}, work, plus \ml{സമയം}, time. You find it on the board at every office door.

\begin{cousinweb}
\textbf{\ml{പ്രവൃത്തിസമയം}} (\emph{pravṛttisamayaṁ}) is a Malayalam compound built from two words borrowed from Sanskrit, \emph{pravṛtti}, \textquotedblleft{}moving onwards, activity\textquotedblright{}, and \emph{samaya}, \textquotedblleft{}time\textquotedblright{} (Monier-Williams).
\end{cousinweb}

\begin{grammarlens}[title={What you've built}]
One more word for getting around.
\end{grammarlens}

\subsection*{Your turn}
\begin{itemize}
\raggedright
  \item \emph{Say it:} \emph{pravṛttisamayaṁ}
  \item \emph{Say it:} \emph{pravṛttisamayaṁ}, once more
  \item \emph{Say it:} \emph{bāṅkinṟe pravṛttisamayaṁ rāvile pattŭ mutal nālŭ vareyāṇŭ.} — The bank's working hours are from ten in the morning to four.
  \item \emph{Recall:} say the Malayalam for a department, then the Malayalam for a branch office, then say \emph{pravṛttisamayaṁ} again
\end{itemize}

\begin{tcolorbox}[breakable,colback=teal!4,colframe=teal!35!black,title={Before you move on}]
What does \ml{പ്രവൃത്തിസമയം} mean? (\textquotedblleft{}Office hours; working hours\textquotedblright{}.) Say it once more.
\end{tcolorbox}
\section[\texorpdfstring{\ml{പ്രവൃത്തിദിനം} (pravṛttidinaṁ)}{പ്രവൃത്തിദിനം (pravṛttidinaṁ)}]{\ml{പ്രവൃത്തിദിനം} (pravṛttidinaṁ) — a working day}

\label{lesson:ML-C486-pravrttidinam}



\begin{quote}
Before the new one: say the Malayalam for a branch office, then the Malayalam for office hours.
\end{quote}

\subsection*{You'll want to know: \ml{പ്രവൃത്തിദിനം}}
\textbf{\ml{പ്രവൃത്തിദിനം}} — \emph{pravṛttidinaṁ} — \textquotedblleft{}a working day\textquotedblright{}.

\begin{quote}
\ml{ശനിയാഴ്ച} \ml{ഇവിടെ} \ml{പ്രവൃത്തിദിനമല്ല}.
\end{quote}

\begin{quote}
\emph{śaniyāḻca iviṭe pravṛttidinamalla.}
\end{quote}

\begin{quote}
Saturday is not a working day here.
\end{quote}

Its opposite is an \ml{അവധിദിനം}, a holiday.

\begin{cousinweb}
\textbf{\ml{പ്രവൃത്തിദിനം}} (\emph{pravṛttidinaṁ}) is a Malayalam compound built from two words borrowed from Sanskrit, \emph{pravṛtti}, \textquotedblleft{}activity\textquotedblright{}, and \emph{dina}, \textquotedblleft{}a day\textquotedblright{} (Monier-Williams).
\end{cousinweb}

\begin{grammarlens}[title={What you've built}]
One more word for getting around.
\end{grammarlens}

\subsection*{Your turn}
\begin{itemize}
\raggedright
  \item \emph{Say it:} \emph{pravṛttidinaṁ}
  \item \emph{Say it:} \emph{pravṛttidinaṁ}, once more
  \item \emph{Say it:} \emph{śaniyāḻca iviṭe pravṛttidinamalla.} — Saturday is not a working day here.
  \item \emph{Recall:} say the Malayalam for a branch office, then the Malayalam for office hours, then say \emph{pravṛttidinaṁ} again
\end{itemize}

\begin{tcolorbox}[breakable,colback=teal!4,colframe=teal!35!black,title={Before you move on}]
What does \ml{പ്രവൃത്തിദിനം} mean? (\textquotedblleft{}A working day\textquotedblright{}.) Say it once more.
\end{tcolorbox}
